% !TeX root = _main.tex
% ========================================
% 結論・今後の課題
% ========================================

%1
\section{おわりに}

本研究では，物語表現という形式を通じて，音による情報伝達の可能性を検討する目的で，音のみを用いた物語作品を制作した．具体的には，音声素材に対するピッチやタイミングなどの演奏情報の制御を行うとともに，ボリュームの調整およびエフェクトの付加を通じて，物語における物理的な動作や，登場キャラクターの心理的な描写の表現を試みた．

作品制作において，筆者自身は，各音声がどのような出来事や感情を示し得るのかを想定して音声設計を行った．しかし，その実演を通じて，音情報のみでは各イベントの内容の把握が難しいという意見が得られた．これは，制作時に想定していた音声の知覚と，聴取時に得られる受け取られ方とが，必ずしも一致しない可能性があることを示している．

以上の内容を踏まえると，音のみを用いた物語表現においては，音声設計の妥当性を検討する上で，比較的小さなまとまりの音声イベントを単位として整理し，その役割を明確化することが重要であると考えられる．本作品は，音のみで物語を構成する試みの初期的な実践として位置づけられ，今後の音響表現による出来事や感情の伝達を検討するための基礎的な知見を提供するものである．